% LaTeX version of the official Apart Research sprint report template
% (Google Doc: https://docs.google.com/document/d/1Y5wd3VRbXU5N08BvAGiJUQTaWbMCSPjhzXsR-ozuehg).
% Layout follows the .docx: US Letter, 1in margins, Old Standard TT, 10pt body,
% 20pt bold title between rules, 14pt bold numbered section headings.
%
% Build: latexmk -pdf main.tex   (or: pdflatex main && bibtex main && pdflatex main && pdflatex main)
%
% Before submitting: set \showguidancefalse to drop every guidance box and
% guidance paragraph, and delete any leftover placeholder text.

\documentclass[10pt]{article}

% ---------- page and fonts ----------
\usepackage[letterpaper,margin=1in]{geometry}
\usepackage[T1]{fontenc}
\usepackage[utf8]{inputenc}
\usepackage{OldStandard}          % Old Standard TT, as in the template
\usepackage[scaled=0.92]{helvet}  % sans font for the guidance box (Arial in the .docx)
\usepackage{microtype}
\usepackage{setspace}
\setstretch{1.15}                 % the template's 1.15 line spacing

% ---------- content packages ----------
\usepackage{amsmath,amssymb}
\usepackage{graphicx}
\usepackage{booktabs}
\usepackage{tabularx}
\usepackage{multirow}
\usepackage[table]{xcolor}
\usepackage{enumitem}
\usepackage{caption}
\usepackage{subcaption}
\usepackage[numbers,sort&compress]{natbib}
\usepackage{tcolorbox}
\usepackage{titlesec}
\usepackage{ulem}\normalem
\usepackage{hyperref}

\definecolor{linkblue}{HTML}{1155CC}
\definecolor{guidancefill}{HTML}{FFF2CC}
\definecolor{subheadgray}{HTML}{434343}
\hypersetup{colorlinks=true,urlcolor=linkblue,linkcolor=linkblue,citecolor=linkblue}

\captionsetup{font=small,labelfont=bf}
\setlist[enumerate]{itemsep=0pt,topsep=4pt}
\setlist[itemize]{itemsep=0pt,topsep=4pt}
\setlength{\parindent}{0pt}
\setlength{\parskip}{6pt}

% ---------- headings (numbered "1. Introduction", bold, 14pt) ----------
\titleformat{\section}{\fontsize{14}{17}\selectfont\bfseries}{\thesection.}{0.4em}{}
\titlespacing*{\section}{0pt}{10pt}{2pt}
\titleformat{\subsection}{\fontsize{12}{15}\selectfont\bfseries}{}{0pt}{}
\titlespacing*{\subsection}{0pt}{8pt}{2pt}
\titleformat{\subsubsection}{\normalsize\bfseries\color{subheadgray}}{}{0pt}{}

% ---------- guidance text (remove before submitting) ----------
\newif\ifshowguidance
\showguidancetrue   % <-- set to \showguidancefalse for the final PDF
\newcommand{\guidance}[1]{\ifshowguidance{#1\par}\fi}
\newcommand{\guidancebox}[1]{\ifshowguidance
  \begin{tcolorbox}[colback=guidancefill,colframe=black,boxrule=0.5pt,arc=0pt,
    left=6pt,right=6pt,top=4pt,bottom=4pt,fontupper=\sffamily\small]#1\end{tcolorbox}\fi}

% ---------- metadata ----------
\newcommand{\projecttitle}{PROJECT TITLE}
\newcommand{\sprintfootnote}{Research conducted at the
  \href{https://apartresearch.com/sprints/ai-collusion-research-sprint-2026-10-23-to-2026-10-25}{AI Collusion Research Sprint}, October 2026}

\begin{document}
\thispagestyle{plain}

% ================= title block =================
\noindent\rule{\linewidth}{3pt}
\vspace{6pt}
\begin{center}
  {\fontsize{20}{24}\selectfont\bfseries \projecttitle\footnote{\sprintfootnote}\par}
\end{center}
\vspace{2pt}
\noindent\rule{\linewidth}{0.8pt}

\vspace{8pt}
\begin{center}
\begin{tabular}{ccc}
  Author name 1 & Author name 2 & Author name 3 \\
  Affiliation   & Affiliation   & Affiliation   \\[8pt]
  Author name 4 & Author name 5 & Author name 6 \\
  Affiliation   & Affiliation   & Affiliation   \\
\end{tabular}

\vspace{8pt}
\textbf{With}\\
Apart Research
\end{center}

% ================= abstract =================
\begin{center}
  {\fontsize{14}{17}\selectfont\bfseries Abstract}
\end{center}
\begin{list}{}{\leftmargin=0.75in\rightmargin=0.75in}\item[]
\itshape
Summarize your project in 150--250 words. A strong abstract lets a reviewer understand
what you did and why it matters without reading anything else.
\uline{Make sure to cover: the problem, your approach, key results, and the main takeaway}.
Polish it last: the abstract should reflect your final results, not your initial plan.
\end{list}

\clearpage

% ================= guidance box =================
\guidancebox{%
\textbf{How to use this template}: Replace the \textit{italicized guidance text} under each
section with your content. The section structure is strong guidance but not rigid. If your
project requires a different organization, feel free to adapt.
\uline{Delete all guidance text including this info box before submitting}. Make sure the
project title and author information above are up to date.

\medskip
\textbf{How your submission will be evaluated}: Your project will be judged on the quality of
this written report. To understand what we look for, see our
\href{https://apartresearch.notion.site/sprint-evaluation-rubric}{Evaluation Rubric}.

\medskip
\textbf{Recommended length}: 4 pages excluding references and appendix.\\
Rough guide: Intro \& Related Work 1p, Methods and Results: 2.5p, Discussion 0.5p.
}

% ================= 1. Introduction =================
\section{Introduction}

\guidance{What problem are you addressing and why does it matter?
Connect it to your work: we want to know why your work is practically valuable.}

\guidance{Provide enough background for readers to understand your work.
If relevant, briefly describe the threat model or failure mode you're addressing; reference
prior work that motivates why it is worth addressing, or explain it yourself.}

\guidance{Aspire to clearly list your most important contributions that go beyond what exists today.}

\textbf{Our main contributions are:}
\begin{enumerate}
  \item {[First contribution --- what new thing did you create, discover, or demonstrate?]}
  \item {[Second contribution]}
  \item {[Third contribution, if applicable]}
\end{enumerate}

% ================= 2. Related Work =================
\section{Related Work}

\guidance{What prior work is most similar, and how does your work differ?
Cite the most relevant papers, tools, or projects. Explain what gap your work addresses.}

\guidance{Some questions which may help:
\begin{enumerate}[label=(\alph*)]
  \item When and why would someone use your method over the existing state-of-the-art?
  \item What information/insight does your method provide which we did not have before?
\end{enumerate}}

% ================= 3. Methods =================
\section{Methods}

\guidance{Describe your approach clearly enough that someone could replicate it. Include key
design choices and justify them where relevant (hint: the more you can back up your design
choices by referencing prior work, the better).}

\guidance{What models/datasets/tools did you use and why? What were key parameters or design
decisions? What did you try that didn't work? Could someone reproduce your work from this
description?}

% ================= 4. Results =================
\section{Results}

\guidance{Present your main findings with appropriate evidence. Use figures and tables where
appropriate (we strongly encourage at least one figure, see tips below). Distinguish between
observations and interpretations.}

\guidance{Argue why your claims are robust. E.g., if your approach ``performs better'' than alternatives:
\begin{enumerate}[label=(\alph*)]
  \item Do you have enough data? Is the difference statistically significant?
  \item Is it robust? Or do small changes to your setup cause substantial changes to your results?
\end{enumerate}}

\guidance{\textbf{Tips for figures and tables:}
\begin{itemize}
  \item Number all figures (Figure 1, Figure 2\dots) and tables (Table 1, Table 2\dots)
  \item Include descriptive captions that can be understood without the main text
  \item Place figures/tables near where they're first referenced
  \item \textbf{Ensure text in figures is legible!}
\end{itemize}}

% Example figure and table (delete or replace):
% \begin{figure}[t]
%   \centering
%   \includegraphics[width=\linewidth]{figures/fig1.pdf}
%   \caption{A caption that can be understood without the main text.}
%   \label{fig:main}
% \end{figure}
%
% \begin{table}[t]
%   \centering
%   \caption{A descriptive caption.}
%   \label{tab:main}
%   \begin{tabular}{lrr}
%     \toprule
%     Condition & $n$ & Rate [95\% CI] \\
%     \midrule
%     A & 0 & 0.00 [0.00, 0.00] \\
%     \bottomrule
%   \end{tabular}
% \end{table}

% ================= 5. Discussion and Limitations =================
\section{Discussion and Limitations}

\guidance{Discuss the broader implications for AI safety.
What do your results mean? What trends do you notice and what might they indicate?}

\subsection{Limitations}

\guidance{What are the limitations of your work? What threat models or failure modes did you
not address? Be honest about constraints --- methodological limitations, scope limitations, or
aspects you couldn't fully address in the hackathon timeframe. Explicitly note the assumptions
you made, whether implicitly or explicitly, and how the interpretation of your results would
change if a given assumption did not hold.}

\subsection{Future Work}

\guidance{What are the natural next steps? How could this work be extended?}

% ================= 6. Conclusion =================
\section{Conclusion}

\guidance{Briefly summarize your main findings and their implications (1--2 paragraphs).}

% ================= unnumbered end matter =================
\section*{Code and Data}

\guidance{Include links if applicable. If your project doesn't involve code (e.g., policy
analysis) or if there are info-hazard considerations, note that here.}

\begin{itemize}
  \item \textbf{Code repository}: [Link to GitHub/GitLab if applicable]
  \item \textbf{Data/Datasets}: [Link if applicable]
  \item \textbf{Other artifacts} (optional): [Demo link, video walkthrough, Hugging Face Space, etc.]
\end{itemize}

\section*{Author Contributions (optional)}

[e.g., ``A.B.\ led the project and designed experiments. C.D.\ implemented the code. All
authors contributed to writing and reviewed the final manuscript.'']

\section*{References}
\guidance{Use a consistent citation format. Include: Author(s), Year, Title, Venue/Publisher,
and URL or DOI where available.}
% Entries live in references.bib; cite with \citep{key} / \citet{key}.
{\renewcommand{\section}[2]{}%  suppress the bibliography's own heading
\nocite{apart2026collusionsprint}% remove once the paper cites real references
\bibliographystyle{plainnat}
\bibliography{references}}

% ================= appendix =================
\appendix
\renewcommand{\thesection}{\Alph{section}}

\section*{Appendix (optional)}
\guidance{Supplementary material such as additional figures, detailed methodology, prompts
used, extended results, etc.}

% Required by the AI Collusion Research Sprint submission guidelines (sprint page,
% Guidelines tab); not part of the generic .docx template.
\section*{Limitations and Dual-Use / Ethical Considerations}
\guidance{Required for this sprint. State what the result does and does not establish, and
any way the setup or findings could make real-world collusion easier to carry out or harder to
detect.}

\section*{LLM Usage Statement}

\guidance{If you used LLM assistance in developing your project or writing this report,
briefly note how. Ensure all claims and results have been verified.

NOTE: We strongly encourage that the final version of the submission is primarily written by your team.}

[e.g., ``We used Claude to brainstorm approaches and help draft sections. All results and
claims were independently verified.'']

\end{document}
